\chapter{Auswertung}
\label{ch:Auswertung}
% Liste der genutzer Formeln für die Fehlerrechnung
\section*{Fehlerrechnung}
Für die statistische Auswertung von $n$ Messwerten $x_i$ werden folgende Größen definiert \cite{errorSkript25}:
\begin{align}
    \bar{x} &= \frac{1}{n} \sum_{i=1}^{n} x_i \vphantom{\sqrt{\sum_i^n}^2} && \text{\textcolor{gray}{Arithmetisches Mittel}} \label{eq:arithmetisches_mittel} \\
    \sigma^2 &= \frac{1}{n-1} \sum_{i=1}^{n} (x_i - \bar{x})^2 \vphantom{\sqrt{\sum_i^n}^2} && \text{\textcolor{gray}{Varianz}} \label{eq:Varianz} \\
    \sigma &= \sqrt{\frac{1}{n-1} \sum_{i=1}^{n} (x_i - \bar{x})^2} \vphantom{\sqrt{\sum_i^n}^2} && \text{\textcolor{gray}{Standardabweichung}} \label{eq:standardabweichung} \\
    \Delta \bar{x} &= \frac{\sigma}{\sqrt{n}} = \sqrt{\frac{1}{n(n-1)} \sum_{i=1}^n(\bar x - x_i)^2} \vphantom{\sqrt{\sum_i^n}^2} && \text{\textcolor{gray}{Fehler des Mittelwerts}} \label{eq:fehler_mittelwert} \\
    \Delta f &= \sqrt{\left(\frac{\partial f}{\partial x} \Delta x\right)^2 + \left(\frac{\partial f}{\partial y} \Delta y\right)^2} \vphantom{\sqrt{\sum_i^n}^2} && \text{\textcolor{gray}{Gauß’sches Fehlerfortpflanzungsgesetz für $f(x,y)$}} \label{eq:gauss_fehlfortpflanzung} \\
    \Delta f &= \sqrt{(\Delta x)^2 + (\Delta y)^2} \vphantom{\sqrt{\sum_i^n}^2} && \text{\textcolor{gray}{Fehler für $f = x + y$}} \label{eq:fehler_summe} \\
    \Delta f &= |a| \Delta x \vphantom{\sqrt{\sum_i^n}^2} && \text{\textcolor{gray}{Fehler für $f = ax$}} \label{eq:fehler_proportional} \\
    \frac{\Delta f}{|f|} &= \sqrt{\left(\frac{\Delta x}{x}\right)^2 + \left(\frac{\Delta y}{y}\right)^2} \vphantom{\sqrt{\sum_i^n}^2} && \text{\textcolor{gray}{relativer Fehler für $f = xy$ oder $f = x/y$}} \label{eq:relativer_fehler} \\
    \sigma &= \frac{|a_{lit} - a_{gem}|}{\sqrt{\Delta a_{lit}^2 + \Delta a_{gem}^2}} \vphantom{\sqrt{\sum_i^n}^2} && \text{\textcolor{gray}{Berechnung der signifikanten Abweichung}} \label{eq:signifikante_abweichung}
\end{align}

\newpage

\section{Absorbiert von Beta--Strahlung in Aluminum}
\label{sec:A1}
Zunächst soll die Absorption von $\beta$--Strahlung durch Aluminum betrachtet werden. Dazu wurde bei verschiedenen Dicken der Aluminumschicht die Zählrate bestimmt. Die Messwerte sind der Tabelle M.3.1 im \hyperref[ch:Protokoll]{Messprotokoll} zu entnemen. Der nullefekt wurde ist in Aufgabe 2 des \hyperref[ch:Protokoll]{Messprotokolles} bestimmt und leigt bei $n_0^\beta = (0.28 \pm 0.03) \mathrm{s^{-1}}$
Die Unegnauigkeit der Zählrate wird nach der Possion--Statistik nach
\eq{possion_err_counts}{
    \Delta n = \frac{\sqrt{n}}{\sqrt{N}}
}
bestimmt. Dabei bezeichnet $N$ die Anzahl der Anzahl der Ereignisse.
Die Messergebnisse sind der \cabb{plots/Absorption_beta_Aluminum} zu entnehmen.
Darüberhinaus wurde eine Regression nach \ceq{zerfallsgesetz} der Form
\eq{reg_expo_decay}{
    N(d) = N_0 \cdot \exp(- \mu \cdot d) - n_0^\beta
}
an die Messdaten gefittet. Die Fitparameter sind der Legende der Abbildung zu netnehmen. 
Es wurde eine Fehlerabschätzung der Regression hinzuefügt. Diese wurde über die in Python bestimmte Ungenauigkeit des Absorptionskoeffizientenes $\mu$ bestimmt, in dem die Regression einmal mit addierten, einmal mit subdrahierten Fehler durhcgeführt wurde. Der entstandene Fehlerschlauch ist in dem Plot rot dargestellt. 

\img[1]{plots/Absorption_beta_Aluminum}{Messdaten der Absorptionsmessung von Beta--Strahlung in Aluminum und Regressionsfunktion.}

Mit den Ergebnissen der Regression können die erwartete maximale Reichweite der $\beta$--Strahlung in Aluminum und die Flächendichte bestimmt werden. 
Die maximale Reichweite ist gerade der Wert, wo die Regression null wird, bzw. sich der Null belibig nah nährt. Hier wurde bestimmt, bei welcher Dicke der Aluminumschicht die Ereignisse pro Sekunde unter $10^{-14}$ geganden sind. Dies ist gerade ungefähr die Genauigkeit von Single--Precision--Floats. Für die Ungenauigkeit wurde das Prozedere für die beiden Fehlerabschätzungsgraphen wiederholt und die mittlere Abweichung nach dem \hyperref[eq:arithmetisches_mittel]{arithmetischen Mittel (\ref*{eq:arithmetisches_mittel})} bestimmt. Es ergibt sich 
\res[-3]{x_\text{max}^\beta}{2.841}{0.017}{m}

Über die maximale Riechweite, der Dichte von Aluminum $\rho_\text{Al} = 2.7 \mathrm{g/cm^3}$ und der Flächendichte $R_\text{ES} = 0.13 \mathrm{g/cm^2}$ kann via
\eq{idszf}{
    R_\beta = \rho_\text{Al} x_\text{max} + R_\text{ES}, \qquad 
    \Delta R_\beta = \rho_\text{Al} Delta x_\text{max}
}
bestimmt werden. Dabei wurde die Ungenauigkeit via \hyperref[eq:gauss_fehlfortpflanzung]{Gauß'scher Fehlerfortpflanzung (\ref*{eq:gauss_fehlfortpflanzung})} bestimmt.  Einsetzen der bekannten Werte liefert

\res{R_\beta}{8.97}{0.05}{\frac{kg}{m^2}}

Über die \cabb{plt1_a} kann die maximale Energie bestimmt werden. Somit liegt die Energie bei
\res{E_\text{exp}^\beta}{1.95}{0.20}{MeV}

Die Ungenauigkeit wurde mit 10\% Unsicherheit geschätzt. Der Literaturwert liegt bei $E_\text{lit}^\beta = 2.274$MeV \cite{skriptPAP22}. Die \hyperref[eq:standardabweichung]{Standardabweichung (\ref*{eq:standardabweichung})} beträgt somit $1.66\sigma$ und ist somit statistisch nicht signifikant.


\section{Absorption von Gamma--Strahlung durch Blei}
\label{sec:A2}
In dieser Sektion wird die Absorption von $\gamma$--Strahlung durch Blei betrachtet. Analgo zu \csec{A1} werden die Ergebnissen in Abhänigkeit der Dicke der Absorptionsschicht betrachtet. Die Messdaten werden geplottet und mit einer Regression nach \ceq{reg_expo_decay} versehen.
Regression, Fitparameter und Messdaten sind der \cabb{plots/Absorption_gamma_Blei} zu entnehmen.
\img[1]{plots/Absorption_gamma_Blei}{Messdaten zur Gamma--Apsorption und Regressionsverlauf.}

In dieser Sektion ist aus der Regression der Parameter
\eq{mu_param}{
    \mu = (0.544 \pm 0.009) \mathrm{cm^{-1}}
}
benötigt. Über diesen kann der Massenschwähungskoeffizient $\zeta = \mu/\rho$. Dessen Ungenauigkeit wird via \hyperref[eq:gauss_fehlfortpflanzung]{Gauß'scher Fehlerfortpflanzung (\ref*{eq:gauss_fehlfortpflanzung})} bestimmt:
\eq{err_zeta}{
    \Delta \zeta = \frac{\Delta \mu}{\rho} \, .
}
Mit der Dichte von Blei $\rho_\mathrm{Pb} = 11.34  g/cm^3$ ergibt sich ein Massenschwähungskoeffizient von
\imp[-3]{\zeta}{4.79}{0.08}{\frac{m^2}{kg}} \label{eq:masskof}
Über diesen Wert kann aus \cabb{plt1_b} die Energie der Photonen zu
\res{E_\text{exp}^\gamma}{1.55}{0.16}{MeV}
bestimmt werden.
Es gibt für $^{60}\mathrm{Co}$ zwei relevante $\gamma$--Übergänge. Es werden hier beide getrennt und einmal kombineirt betrachtet.
Dieser weicht vom Literaturwert $E_\text{lit,1}^\gamma = 1.173$MeV um $2.43\sigma$ ab und vom $E_\text{lit,2}^\gamma = 1.333$MeV um $1.4\sigma$ ab. 
Vom \hyperref[eq:arithmetisches_mittel]{arithmetischen Mittel (\ref*{eq:arithmetisches_mittel})} der beiden Energien um $1.92\sigma$ und somit ist die Abweichung in keinem Fall signifikant. 

\section{Aktivität der Gamma--Strahlung}
\label{sec:A3}
In dieser Sektion wird die Aktivität des Gammastrahlers betrachtet. Dazu werden die Werte der Aufgabe 5 des \hyperref[ch:Protokoll]{Messprotokolls} betrachtet. 
Zunächst soll die Aktivität des Präperates bestimmt werden. 
Die Aktivität wird nach \ceq{aktivitaet_korr} die Aktivität berechnet, wobei hier zunächst keine Korrekturen vorgenommen, werden;
\eq{A0}{
    A_0 = \frac{2 d^{2} n}{\epsilon r^{2}} \, , \qquad
    \Delta A_0 = A_0 \cdot 
            \sqrt{\frac{\Delta n^{2}}{n^{2}} + \frac{4 \Delta d^{2}}{d^{2}}} \, .
}
Dabei beschreibt $d$ den Abstand, $n$ die Zählrate, $\epsilon = 0.04$ ist der konstante Detektoreffizienz nach Praktikumsskript und $r = 0.7$cm den Radius des Zählrohres.
Um den Zerfall in zwei Photonen zu berücksichtigen wurde ebenfalls ein Faktor $1/2$ hinzugefügt. 
Diese, sowie alle weiteren Ergebnisse dieser Aufgabe sind in \ctab{Erg_A3} eingetragen.
Die endliche Länge des Zählrohrs bedingt, dass die Detektoreffizienz auf Teilchen bezogen ist, welche die volle Zählrohrlänge durchlaufen. Ohne entsprechende Korrektur resultiert daraus eine Überschätzung des abgedeckten Raumwinkels und folglich eine zu geringe Bestimmung der Präparataktivität.
Innerhalb der Messdaten zeigt sich ferner, dass die Abweichungen von den Erwartungswerten bei kleineren Abständen ausgeprägter sind. Dies erklärt sich aus der getroffenen Näherung, wonach der Abstand zwischen Zählrohr und Probe wesentlich größer sein soll als der Zählrohrradius -- eine Annahme, deren Gültigkeit mit wachsendem Abstand zunimmt.

Die erste Korrektur berücksichtigt die Länge des Zählrohres $l = 4$cm:
\eq{A1}{
    A_1 = \frac{n \left(2 d + l\right)^{2}}{2 \epsilon r^{2}} \, , \qquad
    \Delta A_1 = A_1 \cdot
            \sqrt{\frac{16 \Delta d^{2}}{\left(2 d + l\right)^{2}} + \frac{\Delta n^{2}}{n^{2}}} \, .
}
Zusätzlich soll der der Korrekturfaktor via 
\eq{korrf}{
    k = \frac{A_1}{A_0} \, , \qquad \Delta k = k \sqrt{\left(\frac{\Delta A_1}{A_1}\right)^2 + \left(\frac{\Delta A_0}{A_0}\right)^2}
}
für alle Messreihen bestimmt werden. 
Aus den Korrekturfaktoren geht hervor, dass die Korrektur der Zählrohrlänge signifikant vom Abstand zwischen Zählrohr und Präparat beeinflusst wird. Dieser Befund stellt einen zweiten Effekt dar, der zur beobachteten Abstandsabhängigkeit der unkorrigierten Aktivitäten beiträgt.

Die zweite Korrektur berücksichtigt weiter die Dichte des Präperates $\rho = 7.9 \mathrm{g/cm^3}$ und die Dicke $b = 0.14\mathrm{cm}$ um die Absorption in der Präperatskapesel zu berücksichtigen. Zudem wird der Massenschwähungskoeffizient nach \ceq{masskof} genutzt. Die Formel lautet somit
\eq{A2}{
    A_2 = \frac{n \left(2 d + l\right)^{2}}{2 \epsilon r^{2}} e^{\zeta \rho b} \, , \qquad
    \Delta A_2 = A_2 \cdot
            \sqrt{\rho^{2} b^{2} \Delta \zeta^{2} + \frac{16 \Delta d^{2}}{\left(2 d + l\right)^{2}} + \frac{\Delta n^{2}}{n^{2}}} \, .
}
Auffällig ist, dass die Absorptionskorrektur die Abweichung der Resultate vom Erwartungswert erhöht. Dies könnte darauf hindeuten, dass die vorangegangene Längenkorrektur bereits zu groß dimensioniert war. Plausibel wird dies durch die Wahl des mittleren Raumwinkels, der auf der Position in der Zählrohramitte beruht, ohne die quadratische Skalierung der Kugeloberfläche mit dem Abstand zu berücksichtigen.

Die theoretische Aktivität berechnet sich via
\eq{A_theo}{
    A_\text{theo} = A_\text{start} \cdot 2^{-T / t_{1/2}} \,.
}
Dabei beschreibt $A_\text{start} = 3700$ kBq die Aktivität bei Herstellung (Präperat SN 372), $T = 5950$ d die Tage zwischen Herstellung des Präperates (02.03.2010) und Versuchsdurchführung (16.06.2026). $t_{1/2} = 1925$ d die Halbwärtszeit von $^{60}$Co.
Somit leigt die erwartete Aktivität bei
\imp[5]{A_\text{theo}}{4.34192}{0}{Bq}
Es wird keine züsätzliche Ungenauigkeit angegeben.

\begin{table}[h!]
    \centering
    \caption{Ergebnisse zur Aufgabe: Aktivität der $\gamma$--Strahlung. $A_\text{theo} = 4.3 \cdot 10^5$ Bq. Abweichungen in $\sigma$.}
    \label{tab:Erg_A3}
    \begin{tabular}{C | C C C | C C C | C C}
        d \; [\mathrm{cm}] & n \; [\mathrm{s^{-1}}] & A_0 \; [\mathrm{10^5 \; Bq}] & \text{Abw. } & A_1 \; [\mathrm{10^5 \; Bq}] & k & \text{Abw. } & A_2 \; [\mathrm{10^5 \; Bq}] & \text{Abw. }\\
        \toprule
         5 & 140.5 \pm 1.6 & 3.6 \pm 1.4 & 0.53 & 7.0 \pm 2.0 & 2.0  \pm 1.0  & 1.33 & 7.4 \pm 2.1 & 1.45 \\
        10 &  40.9 \pm 1.0 & 4.2 \pm 0.8 & 0.2  & 6.0 \pm 1.0 & 1.4  \pm 0.4  & 1.65 & 6.3 \pm 1.1 & 1.87 \\
        20 &  10.0 \pm 0.7 & 4.1 \pm 0.5 & 0.51 & 4.9 \pm 0.6 & 1.21 \pm 0.20 & 1.09 & 5.2 \pm 0.6 & 1.49 \\
        \bottomrule
    \end{tabular}
\end{table}
Es zeigt sich, dass keine der Abweichung statistisch signifikant von dem theoretischen Wert abweichen. 
Interesanterweise jedoch, läss sich der Trend beobachten, dass jede Korrektur zu einer größeren Abweichung führt. Während die unkorrigierten Werte leicht unter dem Erwartungswert liegen, führt jede Korrektur zu einer größeren Überschätzung des Wertes. 
Darüberhinaus scheint es, als würden größere Distanzen zu genuaeren Werten führen, nur die unkorrigierte Aktivität weis hier eine kleine Anomalie auf, hier ist der Wert bei mittlerer Distanz der Akurateste.

\section{Absorption der Alpha--Strahlung in Luft}
\label{sec:A4}
In der letzten Sektion wird die Absorbiert von Alpha--Strahlung in Luft betrachtet. Dafür wird der Luftdruck variiert und jeweils die Ereignisse pro 60 Sekunden gemessen. 
Die \cabb{plots/Absorption_alpha_Luft} zeigt die Messergebnisse. Desweiteren wurde eine lineare Regression nummerisch bestimmt. Die Regressionswerte und Fitparameter sind ebenfalls in der Abbildung zu finden. 
Zuletzt wurde die maximale Zählrate bestimmt und über diese die halbe Zählrate. 

Die lineare Regression hat die Form
\eq{lin_reg}{
    N(p) = mp+b  \,,
}
wobei $p$ der Druck ist, $m$ die Steigung und $b$ der Abzissen--Abschnitt. Die Fitparameter wurden nummerisch bestimmt. In dieser Aufgabe wird lediglich die Steigung 
\eq{val_steig}{
    m = (- 0.82120 \pm 0.00008) \mathrm{hPa \, s} \, 
}
benötig.

\img[1]{plots/Absorption_alpha_Luft}{Messdaten zur Absorption von Alpha--Strahlung in Luft mit Fitparametern und linearer Regression.}

Zudem wird der Druck bei halber Zählrate benötigt. Dieser wird via
\eq{ipsdhjbfg}{
    p_{1/2} = \frac{N_\text{max}/2 - b}{m}\, , \qquad 
    \Delta p_{1/2} = \sqrt{\frac{1}{m^{4}} \left(m^{2} \left(\Delta N_\text{{max}}^{2} + \Delta b^{2}\right) + \Delta m^{2} \left(N_\text{{max}} - b\right)^{2}\right)}
}
berechnet werden, wobei die Ungenauigkeit via \hyperref[eq:gauss_fehlfortpflanzung]{Gauß'scher Fehlerfortpflanzung (\ref*{eq:gauss_fehlfortpflanzung})} bestimmt wurde. Es ergibtsich
\imp[4]{p_{1/2}}{3.03}{0.14}{Pa}

Für die Reichweite der Alpha--Strahlung in Luft ist die Formel
\eq{alpha_reichweit}{
    s_1 = \frac{p \cdot s_{0}}{p_{0}} \, , \qquad
    \Delta s_1 = s_1 \cdot \sqrt{\frac{\Delta s_{0}^{2}}{s_{0}^{2}} + \frac{\Delta p^{2}}{p^{2}}}
}
gegeben. Hier beschreibt $p_0 = 1013$hPa den Normaldruck und $s_0 = (4.45 \pm 0.05)$cm die Distanz zwischen Präperat und Zählrohr.
Es ergibt sich eine Reichweite von 
\imp[-2]{s_1}{1.33}{0.06}{m}
Es werden zwei Korrekturen unternommen:
\eq{korrektur_strecken}{
    s_2 = \frac{\rho_g}{1.43 \;[\mathrm{mg/cm^2}]} \cdot 1 \; [\mathrm{cm}]\, , \quad s_3 = 0.68 \; [\mathrm{cm}] \,.
}
Dabei ist $\rho_g = 2.35 [\mathrm{mg/cm^2}]$ im \hyperref[ch:Protokoll]{Messprotokoll} vermerkt. Die beiden Korrekturen werden als genau angenommen. Die finale Reichweite ist die Summe aller drei strecken $s_i$:
\res[-2]{s}{3.65}{0.06}{m}
Es wird die Energie nach \cabb{plt1_a} für dei Alpha--Strahlung abgelesen und ergibt sich zu 
\res{E_\text{exp}^\alpha}{5.2}{0.5}{MeV}
Zum Literaturwert von $E_\text{lit}^\alpha = 5.48$MeV weicht dieser mit $0.44\sigma$ nicht signifikant ab. 